\section{Conventions and finite-part normalization}
\label{conv:section}

We use $\psi=\Gamma'/\Gamma$, $\psi^{(0)}=\psi$, and
$H_n^{(r)}=\sum_{j=1}^n j^{-r}$, with $H_n=H_n^{(1)}$ and
$H_0^{(r)}=0$. A prime on $\zeta(s)$ denotes differentiation
in its order $s$; a superscript $[k]$ on $\Li_s$ also denotes
order differentiation. Derivatives of $\gamma_m(x)$ are with
respect to its argument $x$.

The generalized Stieltjes convention is
\begin{equation}\label{conv:stieltjes}
\zeta(s,x)=\frac1{s-1}+
\sum_{m\ge0}\frac{(-1)^m\gamma_m(x)}{m!}(s-1)^m.
\end{equation}
Thus $\gamma_0(x)=-\psi(x)$ and
$\gamma_m=\gamma_m(1)$. After setting $s=1-u$, the pole-removed
germ is
\[
\zeta(1-u,x)+\frac1u
=\sum_{m\ge0}\gamma_m(x)\frac{u^m}{m!}.
\]
This positive-sign $u$ expansion must not be confused with the
alternating-sign expansion of $\zeta(1+u)$.

Strict multiple zeta functions have decreasing indices:
\[
Z_r(s_1,\ldots,s_r)=
\sum_{n_1>\cdots>n_r\ge1}\prod_{j=1}^r n_j^{-s_j}.
\]
We also write $\zeta(a,b)=Z_2(a,b)$ when the sum converges.
The coefficients $\eta$ and $\theta$ in Section~\ref{q4:section}
are coefficients of $u^1$, with no hidden factorial or sign.
Bernoulli numbers have $B_1=-1/2$, and Bernoulli polynomials use
\[
\frac{te^{xt}}{e^t-1}=\sum_{n\ge0}B_n(x)\frac{t^n}{n!}.
\]

\begin{definition}[Coordinate finite part]
Suppose an integrand on $(0,1)$ has, to every required finite order,
endpoint expansions in complex powers and powers of logarithms,
with an integrable remainder after finitely many subtractions.
Its coordinate finite part is the constant coefficient of the
cutoff integral in the endpoint coordinates $x$ and $1-x$.
One may introduce independent cutoffs at the two ends; the
constant is the coefficient of both cutoff powers zero and both
logarithm powers zero. A common cutoff gives the same constant
for the separate endpoint expansions used here.
\end{definition}

For example,
\[
\FP\int_0^1x^{-j}\,\dd x=-\frac1{j-1}\quad(j\ge2),
\qquad
\FP\int_0^1\frac{\log^k x}{x}\,\dd x=0\quad(k\ge0).
\]
At a general nonzero exponent $\alpha$ the regularized monomial
integral is $1/\alpha$; its parameter derivatives are the
corresponding logarithmic integrals. At $\alpha=0$ the coordinate
constant is zero, whereas the meromorphic function has a pole.
This elementary difference becomes a finite contact term when a
coefficient also vanishes at the same parameter.

A change of endpoint coordinate can change a finite part. Every
identity in this article fixes the coordinate before taking the
constant, differentiating parameters, or forming a primitive.
In particular, the spectral values $D_{p,q}(-2r)$ are initially
defined by meromorphic continuation, and their equality to
ordinary subtracted sums is proved in
Proposition~\ref{tail:ordinary}; it is not an implicit choice
of summation method.

All primitives are normalized. The quartic primitive has zero
constant term at its singular origin. Primitives in the nonreal
resolvent parameter are differences at a specified base point
within one half-plane. Complex powers of the resolvent use the
continuous logarithm specified in
Section~\ref{power:sec}; no branch is inferred from an
unqualified power of a complex number.

